% =============================================================
% Präambel — Pakete, Schriften, Geometrie, Sprache, Kopf-/Fußzeile.
%
% Wird von buch.tex eingebunden, NACH \documentclass.
% Layout-Bausteine (Aufgaben-Umgebungen) sind in befehle.tex,
% Farben in farben.tex.
% =============================================================

% --- Geometrie nach doku/LAYOUT_BRIEFING.md ---
\usepackage[a4paper,
            inner=2.5cm, outer=4.5cm,
            top=2.0cm, bottom=2.5cm,
            marginparwidth=2.7cm, marginparsep=0.15cm,
            headheight=15pt]{geometry}

% --- Schriften ---
% Source Pro Familie als Hausschriften, vendored unter latex/fonts/
% (siehe latex/fonts/README.md für Provenance). Source Serif für den
% Fließtext, Source Sans für Überschriften/Marginalien, Source Code
% für Quelltexte. Math-Formeln werden weiter unten via unicode-math
% gesetzt (STIX Two Math passt visuell zu Source Serif).
\usepackage{fontspec}
\setmainfont{SourceSerif4}[
  Path=fonts/source-serif/,
  Extension=.otf,
  UprightFont=*-Regular,
  ItalicFont=*-It,
  BoldFont=*-Bold,
  BoldItalicFont=*-BoldIt,
  Scale=1.00,
]
\setsansfont{SourceSans3}[
  Path=fonts/source-sans/,
  Extension=.otf,
  UprightFont=*-Regular,
  ItalicFont=*-It,
  BoldFont=*-Bold,
  BoldItalicFont=*-BoldIt,
  Scale=0.95,
]
\setmonofont{SourceCodePro}[
  Path=fonts/source-code-pro/,
  Extension=.otf,
  UprightFont=*-Regular,
  ItalicFont=*-It,
  BoldFont=*-Bold,
  BoldItalicFont=*-BoldIt,
  Scale=0.88,
]
% Marginalien-Symbole (✏ ⌨ ⚙) brauchen einen Symbol-fähigen Font.
% Source Sans hat die Glyphen nicht; wir behalten DejaVu Sans als
% Fallback nur dafür (System-Font, im Container via fonts-dejavu).
\newfontfamily\markerfont{DejaVu Sans}

% --- Sprache, Mikrotypografie ---
\usepackage[ngerman]{babel}
% babel-ngerman macht ASCII-" zum Shorthand-Aktivator ("u → ü, "f → ff
% etc.). Das hat schon zu absurden Verschmelzungen wie "echterFFehler"
% statt "echter" Fehler" geführt. Wir schalten den Shorthand ab; für
% typografisch korrekte deutsche Anführungszeichen nutzen wir Unicode
% („…") direkt im Quelltext.
\shorthandoff{"}
\usepackage{microtype}

% --- Mathematik (\bmod, \le, \ge, \checkmark, …) ---
% amsmath für Mathematik-Umgebungen (align, gather, …).
% unicode-math + STIX Two Math ersetzt die Computer-Modern-Math-
% Schriften; Stilfamilie passt visuell zu Source Serif.
\usepackage{amsmath}
\usepackage{unicode-math}
\setmathfont{STIXTwoMath-Regular.otf}[
  Path=fonts/stix-two/,
]

% --- Glossar (Polish-Sprint M6) ---
% automake=immediate ruft makeglossaries per \write18 auf (shell-escape
% ist eh schon aktiv wegen minted), nonumberlist unterdrückt die
% Seitenzahlen hinter den Glossar-Einträgen (sonst stehen dort die
% Seiten, an denen \gls genutzt wurde).
\usepackage[automake=immediate, toc, nonumberlist]{glossaries-extra}
\makeglossaries

% --- Sachindex ---
\usepackage{imakeidx}
\makeindex[intoc, columns=2]

% --- Farben & Hyperlinks ---
\usepackage{xcolor}
\usepackage[hidelinks]{hyperref}
% pdfpagelabels ist hyperref-Default; wir setzen nur plainpages=false,
% damit hyperref beim \pagenumbering-Wechsel Frontmatter→Mainmatter
% keine Duplikat-Warnung wirft und PDF-Reader die Buchseiten-Labels
% (i, ii, …, 1, 2, …) korrekt anzeigen.
\hypersetup{plainpages=false}

% --- Layout-Bausteine ---
% tcolorbox: nur die Libraries laden, die wir tatsächlich brauchen
% (skins für `enhanced` und `borderline`, breakable für mehrseitige Boxen).
% [most] würde tikzfill u.a. mitziehen, was wir nicht nutzen.
\usepackage{tcolorbox}
\tcbuselibrary{skins, breakable}
\usepackage{changepage}
% \strictpagecheck macht \checkoddpage zuverlässig auch bei Page-Breaks
% (sonst kann die Seitenparität pro Page-Part falsch geprüft werden).
\strictpagecheck

% --- TikZ (für Venn-Diagramme etc.) ---
\usepackage{tikz}
\usetikzlibrary{positioning, calc}

% --- Tabellen ---
\usepackage{booktabs, array, longtable}

% --- Listen mit (a)/(b)/(c)-Labels für Aufgabenteile ---
\usepackage{enumitem}

% --- needspace: erzwingt Mindestplatz vor Aufgabenblöcken ---
\usepackage{needspace}

% --- Code-Highlighting via minted ---
% minted ruft Pygments per Shell auf; deshalb -shell-escape in .latexmkrc.
\usepackage{minted}
% Hintergrundfarbe für Code-Blöcke (sehr helles Grau, naive CMYK aus
% RGB(246, 247, 248), entspricht der bisherigen pandoc-shadecolor).
\definecolor{codebg}{cmyk}{0.01, 0.00, 0.00, 0.03}
\setminted{
  style=tango,
  fontsize=\small,
  bgcolor=codebg,
  frame=none,
  breaklines=true,
  autogobble=true,
  linenos=false,
  baselinestretch=1.05,
  xleftmargin=2pt,
  xrightmargin=2pt,
}

% --- Kopf-/Fußzeile via scrlayer-scrpage ---
\usepackage[autooneside=false, automark]{scrlayer-scrpage}
\clearpairofpagestyles
\automark[chapter]{chapter}
\ihead{\small\textit{\headmark}}
\ohead{\small\pagemark}
\KOMAoptions{headsepline=0.3pt}

% --- Witwen / Waisen vermeiden ---
\widowpenalty=10000
\clubpenalty=10000

% --- Etwas Toleranz für die schmaleren Spalten ---
\setlength{\emergencystretch}{3em}
